% Original source: Zhou, physical PDF pp. 197--201 (printed pp. 181--185).
% Figures: none.
% Uncertainty log: footnote 11 visibly continues beyond physical page 201.

% ZHOU-SOURCE-PAGE: 197 PRINTED: 181

A more efficient approach uses the divide-and-conquer paradigm. Let's define

\[
T(i)=\sum_{x=1}^{i} A[x]
\]

and $T(0)=0$, then $V(i,j)=T(j)-T(\ii-1)$, $\forall 1\leq i\leq j\leq n$. Clearly for
any fixed $j$, when $T(\ii-1)$ is minimized, $V(i,j)$ is maximized. So the maximum
subarray ending at $j$ is $V_{\max}=T(j)-T_{\min}$ where
$T_{\min}=\min(T(1),\ldots,T(j-1))$. If we keep track of and update $V_{\max}$ and
$T_{\min}$ as $j$ increases, we can develop the following $O(n)$ algorithm:

\begin{lstlisting}
T = A[1]; Vmax = A[1]; Tmin = min(0,T)
For j = 2 to n
{
    T = T + A[j];

    If T-Tmin > Vmax then Vmax = T-Tmin;

    If T < Tmin, then Tmin = T;
}
Return Vmax;
\end{lstlisting}

The following is a corresponding C++ function that returns $V_{\max}$ and indices $i$ and $j$
given an array and its length:

\begin{lstlisting}[language=C++]
double maxSubarray(double A[], int len, int &i, int &j)
{
        double T=A[0], Vmax=A[0];
        double Tmin = min(0.0, T);
        for(int k=1; k<len; ++k)
        {
                T+=A[k];
                if (T-Tmin > Vmax) {Vmax=T-Tmin; j=k;}
                if (T<Tmin) {Tmin = T; i = (k+1<j)? (k+1):j;}
        }
        return Vmax;
}
\end{lstlisting}

Applying it to the following array $A$,

\begin{lstlisting}[language=C++]
double A[]={1.0, 2.0, -5.0, 4.0, -3.0, 2.0, 6.0, -5.0, -1.0};
int i = 0, j =0;
\end{lstlisting}

% ZHOU-SOURCE-PAGE: 198 PRINTED: 182

\begin{lstlisting}[language=C++]
double Vmax = maxSubarray(A, sizeof(a)/sizeof(A[1]), i, j);
\end{lstlisting}

will give $V_{\max}=9$, $i=3$ and $j=6$. So the subarray is $[4.0,-3.0,2.0,6.0]$.\qedhere

\end{solution}

\subsection{The Power of Two}

There are only 10 kinds of people in the world---those who know binary, and those who
don't. If you happen to get this joke, you probably know that computers operate using
the binary (base-2) number system. Instead of decimal digits 0--9, each bit (binary digit)
has only two possible values: 0 and 1. Binary representation of numbers gives some
interesting properties that are widely explored in practice and makes it an interesting
topic to test in interviews.

\subsubsection*{Power of 2}

\begin{problem}{How do you determine whether an integer is a power of 2?}
\end{problem}

\begin{solution}
Any integer $x=2^n$ ($n\geq 0$) has a single bit (the $(n+1)$th bit from the right) set
to 1. For example, 8 ($=2^3$) is expressed as $0\cdots01000$. It is also easy to see that
$2^n-1$ has all the $n$ bits from the right set to 1. For example, 7 is expressed as
$0\cdots00111$. So $2^n$ and $2^n-1$ do not share any common bits. As a result,
$x\ \&\ (x-1)=0$, where $\&$ is a bitwise AND operator, is a simple way to identify
whether the integer $x$ is a power of 2.\qedhere

\end{solution}

\subsubsection*{Multiplication by 7}

\begin{problem}{Give a fast way to multiply an integer by 7 without using the multiplication (*) operator?}
\end{problem}

\begin{solution}
\texttt{(x << 3) - x}, where \texttt{<<} is the bit-shift left operator. \texttt{x << 3} is
equivalent to \texttt{x*8}. Hence \texttt{(x << 3) - x} is \texttt{x*7}.\footnote[8]{The
result could be wrong if \texttt{<<} causes an overflow.}\qedhere

\end{solution}

\subsubsection*{Probability simulation}

\begin{problem}{You are given a fair coin. Can you design a simple game using the fair coin so that your
probability of winning is $p$, $0<p<1$?\footnote[9]{Hint: Computer stores binary numbers
instead of decimal ones; each digit in a binary number can be simulated using a fair coin.}}
\end{problem}

% ZHOU-SOURCE-PAGE: 199 PRINTED: 183

\begin{solution}
The key to this problem is to realize that $p\in(0,1)$ can also be expressed as a binary
number and each digit of the binary number can be simulated using a fair coin. First, we
can express the probability $p$ as binary number:

\[
p=0.p_1p_2\cdots p_n=p_1 2^{-1}+p_2 2^{-2}+\cdots+p_n 2^{-n},
\qquad p_i\in\{0,1\},\ \forall i=1,2,\ldots,n.
\]

Then, we can start tossing the fair coin, and count heads as 1 and tails as 0. Let
$s_i\in\{0,1\}$ be the result of the $i$-th toss starting from $i=1$. After each toss, we
compare $p_i$ with $s_i$. If $s_i<p_i$, we win and the coin tossing stops. If $s_i>p_i$,
we lose and the coin tossing stops. If $s_i=p_i$, we continue to toss more coins. Some
$p$ values (e.g., $1/3$) are infinite series when expressed as a binary number
($n\to\infty$). In these cases, the probability to reach $s_i\neq p_i$ is 1 as $i$ increases.
If the sequence is finite, (e.g., $1/4=0.01$) and we reach the final stage with $s_n=p_n$,
we lose (e.g., for $1/4$, only the sequence 00 will be classified as a win; all other three
sequences 01, 10 and 11 are classified as a loss). Such a simulation will give us
probability $p$ of winning.\qedhere

\end{solution}

\subsubsection*{Poisonous wine}

\begin{problem}{You've got 1000 bottles of wines for a birthday party. Twenty hours before the party,
the winery sent you an urgent message that one bottle of wine was poisoned. You happen
to have 10 lab mice that can be used to test whether a bottle of wine is poisonous. The
poison is so strong that any amount will kill a mouse in exactly 18 hours. But before the
death on the 18th hour, there are no other symptoms. Is there a sure way that you can find
the poisoned bottle using the 10 mice before the party?}
\end{problem}

\begin{solution}
If the mice can be tested sequentially to eliminate half of the bottles each time, the
problem becomes a simple binary search problem. Ten mice can identify the poisonous
bottle in up to 1024 bottles of wines. Unfortunately, since the symptom won't show up
until 18 hours later and we only have 20 hours, we cannot sequentially test the mice.
Nevertheless, the binary search idea still applies. All integers between 1 and 1000
can be expressed in 10-bit binary format. For example, bottle 1000 can be labeled as
1111101000 since $1000=2^9+2^8+2^7+2^6+2^5+2^3$.

Now let mouse 1 take a sip from every bottle that has a 1 in the first bit (the lowest bit
on the right); let mouse 2 take a sip from every bottle with a 1 in the second bit; $\ldots$;
and, finally, let mouse 10 take a sip from every bottle with a 1 in the 10th bit (the highest
bit). Eighteen hours later, if we line up the mice from the highest to the lowest bit and
treat a live mouse as 0 and a dead mouse as 1, we can easily back track the label of the
poisonous bottle. For example, if the 6th, 7th, and 9th mice are dead and all others are
alive, the line-up gives the sequence 0101100000 and the label for the poisonous bottle
is $2^8+2^6+2^5=352$.

\end{solution}

% ZHOU-SOURCE-PAGE: 200 PRINTED: 184

\subsection{Numerical Methods}
\enlargethispage{1pt}

The prices of many financial instruments do not have closed-form analytical solutions.
The valuation of these financial instruments relies on a variety of numerical methods. In
this section, we discuss the application of Monte Carlo simulation and finite difference
methods.

\subsubsection*{Monte Carlo simulation}

Monte Carlo simulation is a method for iteratively evaluating a deterministic model using
random numbers with appropriate probabilities as inputs. For derivative pricing, it
simulates a large number of price paths of the underlying assets with probability
corresponding to the underlying stochastic process (usually under risk-neutral measure),
calculates the discounted payoff of the derivative for each path, and averages the
discounted payoffs to yield the derivative price. The validity of Monte Carlo simulation
relies on the law of large numbers.

Monte-Carlo simulation can be used to estimate derivative prices if the payoffs only depend
on the final values of the underlying assets, and it can be adapted to estimate prices if
the payoffs are path-dependent as well. Nevertheless, it cannot be directly applied to
American options or any other derivatives with early exercise options.

\begin{problem}{A. Explain how you can use Monte Carlo simulation to price a European call option?}
\end{problem}

\begin{solution}
If we assume that stock price follows a geometric Brownian motion, we can simulate
possible stock price paths. We can split the time between $t$ and $T$ into $N$ equally-spaced
time steps.\footnote[10]{For European options, we can simply set $N=1$. But for more general
options, especially the path-dependent ones, we want to have small time steps and therefore
$N$ should be large.} So $\Delta t=\dfrac{T-t}{N}$ and $t_i=t+\Delta t\cdot \ii$, for
$i=0,1,2,\ldots,N$.

\Needspace{6\baselineskip}
We then simulate the stock price paths under risk-neutral probability
using equation

\[
S_i=S_{i-1}e^{(r-\sigma^2/2)(\Delta t)+\sigma\sqrt{\Delta t}\epsilon_i},
\]

where $\epsilon_i$'s are IID random variables from standard normal distribution. Let's say
that we simulate $M$ paths and each one yields a stock price $S_{T,k}$, where
$k=1,2,\ldots,M$, at maturity date $T$.

% ZHOU-SOURCE-PAGE: 201 PRINTED: 185

The estimated price of the European call is the present value of the expected payoff,
which can be calculated as

\[
C=e^{-r(T-t)}\frac{\displaystyle\sum_{k=1}^{M}\max(S_{T,k}-K,0)}{M}.\qedhere
\]

\end{solution}

\begin{problem}{B. How do you generate random variables that follow $N(\mu,\sigma^2)$ (normal
distribution with mean $\mu$ and variance $\sigma^2$) if your computer can only generate
random variables that follow continuous uniform distribution between 0 and 1?}
\end{problem}

\begin{solution}
This is a great question to test the basic knowledge of random number generation, the
foundation of Monte Carlo simulation. The solution to this question can be dissected to two
steps.

\begin{enumerate}
\item Generate random variable of $x\sim N(0,1)$ from uniform random number generator
using inverse transform method and rejection method.
\item Scale $x$ to $\mu+\sigma x$ to generate the final random variables that follow
$N(\mu,\sigma^2)$.
\end{enumerate}

The second step is straightforward; the first step deserves some explanations. A popular
approach to generating random variables is the inverse transform method: For any continuous
random variable $X$ with cumulative density function $F$ ($U=F(X)$), the random variable
$X$ can be defined as the inverse function of $U$: $X=F^{-1}(U)$, $0\leq U\leq 1$.
It is obvious that $X=F^{-1}(U)$ is a one-to-one function with $0\leq U\leq 1$. So any
continuous random variable can be generated using the following process:

\begin{itemize}
\item Generate a random number $u$ from the standard uniform distribution.
\item Compute the value $x$ such that $u=F(x)$ as the random number from the distribution
described by $F$.
\end{itemize}

For this model to work, $F^{-1}(U)$ must be computable. For standard normal distribution,

\[
U=F(X)=\int_{-\infty}^{x}\frac{1}{\sqrt{2\pi}}e^{-x^2/2}\,\dd x.
\]

\Needspace{5\baselineskip}
The inverse function has no analytical solution. Theoretically, we can come up with the
one-to-one mapping of $X$ to $U$ as the numeric solution of ordinary differential equation
$F'(x)=f(x)=\dfrac{1}{\sqrt{2\pi}}e^{-x^2/2}$ using numerical integration method such as
the Euler method.\footnotemark[11] Yet this approach is less efficient than the rejection
method:
